\documentclass[11pt,a4paper]{article}

\usepackage[T1]{fontenc}
\usepackage[utf8]{inputenc}
\usepackage{lmodern}
\usepackage[margin=25mm]{geometry}
\usepackage{microtype}
\usepackage{amsmath}
\usepackage{hyperref}

\hypersetup{
    colorlinks=true,
    linkcolor=blue,
    urlcolor=blue,
    pdftitle={An Archival Investigation of J1237+3314},
    pdfauthor={Yashaswy Akella}
}

\setlength{\parindent}{0pt}
\setlength{\parskip}{0.6em}

\newcommand{\arcsec}{\ensuremath{^{\prime\prime}}}
\newcommand{\OII}{\ensuremath{[\mathrm{O\,II}]}}

\title{
    \textbf{What Is Beside This Quasar?}\\[0.4em]
    \large An Archival Investigation of J1237+3314
}

\author{
    Yashaswy Akella\\
    \small Independent astronomy project
}

\date{\today}

\begin{document}

\maketitle

\begin{quote}
\textbf{Research status.}
This is an exploratory project report, not a peer-reviewed
publication. The central measurements have not yet been
independently reproduced. The project does not establish a
new gravitational lens or a confirmed dual active galactic nucleus.
\end{quote}

\section{The question}

A known quasar has a faint neighbour in the sky.

A quasar is a very bright galactic nucleus powered by matter
accreting onto a supermassive black hole.

Two nearby points of light might be separate, physically
associated objects. They might instead be unrelated objects
at different distances, or multiple images of one distant
source produced by gravitational lensing.

This project began as a search for lensed quasars and developed
into a closer examination of the already known pair J1237+3314.

\section{The two components}

Component A is the bright quasar SDSS J123704.05+331423.8,
at redshift approximately 1.288.

Component B is a fainter optical source about \(1.62\arcsec\)
from A. The precise position of the compact radio source
JVAS B1234+335 lies very close to B.

Their Gaia DR3 identifiers are:
\begin{itemize}
    \item A: \texttt{1514739204708610432}.
    \item B: \texttt{1514739204711602816}.
\end{itemize}

The objects and candidate pair were already catalogued.
This project does not claim their discovery.

\section{What data were examined?}

The investigation combines optical positions, radio
astrometry and imaging, optical images, and spectroscopy
from public astronomical archives.

The data include Gaia, DESI, the DESI Legacy Imaging Surveys,
JVAS, FIRST, VLASS, and the Radio Fundamental Catalog.

These observations have different resolutions, apertures,
calibrations, and observing dates. A measurement near the pair
cannot automatically be assigned to one component.

\section{What the analysis reports}

The numerical values below are outputs reported by the project.
They have not been independently verified by the preparation
of this explanatory report.

\subsection{The radio position agrees with B}

The radio position is approximately 4.7 milliarcseconds from
B's Gaia position and approximately \(1.619\arcsec\) from A.

The reported agreement is consistent with the combined
positional uncertainties. It supports associating compact
radio emission with B's position.

It does not establish B's distance or prove that it is
a second nucleus physically associated with A.

\subsection{A spectral feature appears towards B}

The DESI spectrum collected through the fibre centred on B
contains an emission feature with the reported properties:

\begin{itemize}
    \item Wavelength:
    \(8528.6 \pm 0.4\)~\AA.

    \item Integrated flux:
    \((2.78 \pm 0.16)\times10^{-16}\)
    \(\mathrm{erg\,s^{-1}\,cm^{-2}}\).

    \item Intrinsic full width at half maximum under a
    single-Gaussian model:
    \(550 \pm 30\ \mathrm{km\,s^{-1}}\).
\end{itemize}

The quoted errors are statistical fitting estimates, not a
complete account of systematic uncertainties. The flux follows
the DESI point-source calibration.

The feature is reported in both individual exposures and
in an additional extraction from the preprocessed detector
images. It is present before correction for contamination
from A.

These are internal checks of the same underlying observations,
not independent observational confirmation.

\subsection{The feature could be oxygen emission}

One possible identification is the \OII\ doublet, emitted
by singly ionized oxygen.

If correct, the gas redshift is approximately
\(1.2872\)--\(1.2879\), close to A's redshift.

The profile does not uniquely establish the doublet.
Tests of other expected lines provide evidence against
some alternative identifications, but other possibilities remain.

Even if the oxygen identification is correct, the gas could
belong to the surroundings of A rather than to B.

\subsection{Nearby light complicates the spectra}

Atmospheric blurring sends light from A into the fibre
pointed towards B.

The narrow feature persists over the tested contamination
corrections. A corresponding significant feature is not found
in the spectrum centred on A.

However, the two spectra were collected at different dates
and under different seeing conditions. Their comparison does
not precisely locate the gas.

\section{What remains unresolved?}

The project has not established B's nuclear redshift
or confirmed a dual active galactic nucleus.

Possible explanations include:
\begin{itemize}
    \item A distinct active nucleus physically associated with A.
    \item Gas around A extending towards B.
    \item Gas related to an interaction.
    \item An unrelated radio source projected onto nearby gas.
\end{itemize}

A compact radio feature connected to a jet from A is another
alternative considered in the technical investigation.

The optical and radio brightness patterns motivated a lensing
test, but dust, microlensing, source structure, and aperture
differences complicate that comparison. Gravitational lensing
has not been definitively excluded.

Radio sensitivity thresholds and spatial constraints in the
technical analysis are model-dependent supporting results.

\section{Verification and reproducibility}

The repository should contain one matching release of the
scripts, input references, outputs, and reports.

The main documentation files are:
\begin{itemize}
    \item \path{release/README.md}.
    \item \path{release/data_provenance.md}.
    \item \path{release/out/results_table.md}.
\end{itemize}

A successful execution shows that a workflow runs.
It does not automatically validate its assumptions,
uncertainty estimates, or physical interpretation.

Repeated runs by the author demonstrate repeatability.
Independent reproduction is a separate step.

If execution fails, the error and software environment
should be documented rather than hidden.

\section{The most useful next step}

The immediate priority is independent examination of the
spectral feature and radio--optical positional association.

Spatially resolved spectroscopy could test where the gas lies
and identify additional lines. Higher-resolution radio imaging
could clarify radio emission and source structure.

These observations address different questions and would not
automatically confirm a dual AGN.

\section{Conclusion}

This project re-examines an already known ambiguous pair.
It reports a close radio--optical positional association
and a spectral feature that may help interpret the system.

The physical relationship remains unresolved.

The purpose of sharing the work is to document the methods,
reported evidence, and uncertainties---not to announce
a confirmed discovery.

\section*{Data credits and AI assistance}

Original survey observations retain their providers' terms
and acknowledgement requirements. Detailed references and
input provenance are recorded in the technical release.

This explanatory report was prepared
with assistance from ChatGPT (OpenAI).

\section*{Selected data resources}

\begin{itemize}
    \item Gaia:
    \url{https://gea.esac.esa.int/archive/}

    \item DESI public data:
    \url{https://data.desi.lbl.gov/public/dr1/}

    \item Legacy Surveys:
    \url{https://www.legacysurvey.org/}

    \item SIMBAD:
    \url{https://simbad.cds.unistra.fr/}

    \item VizieR:
    \url{https://vizier.cds.unistra.fr/}
\end{itemize}

These links identify data services. They do not replace
the detailed citations and input references in the
technical documentation.

\end{document}